\scanpage{74}
\mspage{\uline{16}}
Relations avec les pts \deleted{\unsure{d'ordre|d'}} rationnels \inserted{et la cohomologie} des % 1
Variétés abéliennes sur un \deleted{corps de fonctions} \inserted{schéma de type fini} \ldots % 2

Soit $A_K$ \unsure{VA|NA} sur corps $K$ de \uline{type fini} % 3
pour commencer. Alors $K$ est le corps des % 4
fonctions d'un schéma $\deleted{\unsure{X}}S$ irréductible et % 5
régulier, de type fini sur $\mathbb{Z}$, et \inserted{tel que} $A_K$ % 6
provient d'un schéma abélien sur $\deleted{\unsure{X}}S$. % 7

De plus, on a % 8
\[
 A_K(K)=A(\deleted{\unsure{X}}S) % 9
\]
et on sait que c'est un \uline{groupe de type % 10
fini} par Mordell-Weil-Néron-Lang. % 11

Soit $f:A\longrightarrow S$ la projection, et % 12
considérant le motif $\mathbb{R}^1 f^{*}_*(\mathbb{Q}_{A^*}(\deleted{\unsure{0}}1))\simeq
\deleted{\unsure{H_1(A)_{\mathbb{Q}}}} \inserted{T_{\text{\unsure{motif|mot}}}(A)}$ % 13
\inserted{[où $A^*$ est le dual de $A$]} % 14
qui est un motif ct tordu sur $S$ \deleted{de} % 15
pur de poids $1$, qui redonne $A_K$ % 16
\deleted{\unsure{à}} isogénie près \deleted{\unsure{sur $K$}} \deleted{\unsure{sur}} (sur $K$ ?) comme % 17
conséquence des conjectures de Tate. % 18
% archive sheet number: 73
